\documentclass{article}
\usepackage[margin=1in]{geometry}
\usepackage{amsmath,amssymb}
\usepackage{standalone}
\usepackage{graphicx}
\usepackage{booktabs}
\usepackage{hyperref}

\title{Mathematical Structure of the 4D GSPT-SBL Governing System:\\
Analysis of Nocturnal Boundary Layer Dynamics over CASES99}

\author{
    \textbf{David E. England, PhD} \\
    Research Engineer III \\
    University of Alabama in Huntsville (UAH) \\
    Email: \href{mailto:david.england@uah.edu}{david.england@uah.edu} \\
    ORCID: \href{https://orcid.org/0000-0002-1825-0097}{0000-0002-1825-0097}
}
\date{\small Assembled Automatically: July 13, 2026}

\begin{document}

\maketitle

\begin{abstract}
This manuscript details the formulation and geometric structures of a four-dimensional stable boundary layer (SBL) model, analyzed using Geometric Singular Perturbation Theory (GSPT). The dynamic system couples fast turbulent kinetic energy evolution with slow atmospheric rotational and skin temperature variables, validated against the CASES99 field dataset.
\end{abstract}

\section{Introduction and Theoretical Framework}
\label{sec:theory}

The planetary boundary layer under strongly stable stratification exhibits complex, non-equilibrium phase transitions often characterized by sudden turbulence collapses and regime tracking. Traditional continuous-prose derivations frequently obscure the deep fast-slow alignment driving these drops. To resolve this, this framework models the coupled boundary layer system by isolating the short-timescale turbulent dissipation from the longer-timescale macro-environmental forcing.

By applying Geometric Singular Perturbation Theory (GSPT), we project the state trajectories onto a shifting manifold governed by surface radiation and mechanical shear. This maps the mathematical boundaries where turbulence ceases to act as an active continuous background component and collapses down to its regularized residual floor.

\section{The Complete 4D GSPT-SBL Governing System}
\label{sec:governing_system}

The full state vector is defined as $\mathbf{x} = (e, U, V, T_s) \in \mathbb{R}^4$, where $e$ is the fast turbulent coordinate, and $(U, V, T_s)$ are the slow environmental coordinates. The system is scaled using the small singular perturbation parameter $\varepsilon$, representing the ratio of the rapid turbulent allocation timescale to the slow macro-environmental thermodynamic forcing timescale ($0 < \varepsilon \ll 1$).

\subsection{The Complete Fast-Slow Equations}
\label{subsec:fast_slow_eqs}

\subsubsection*{Fast Subsystem (Turbulence Mechanics)}
\begin{equation}
\varepsilon \frac{de}{dt} = \sqrt{e + \delta} \left[ \eta\,\gamma\,(U^2 + V^2) - K \cdot G(T_s) \right] - \frac{(e+\delta)^{3/2}}{l_0} = \mathcal{F}(e,U,V,T_s)
\label{eq:fast_subsystem}
\end{equation}

\subsubsection*{Slow Subsystem (Atmospheric Rotational Dynamics \& Surface Thermodynamics)}
\begin{align}
\frac{dU}{dt} &= f(V - V_g) - \gamma \sqrt{e + \delta} \cdot U \label{eq:slow_u} \\
\frac{dV}{dt} &= -f(U - U_g) - \gamma \sqrt{e + \delta} \cdot V \label{eq:slow_v} \\
\frac{dT_s}{dt} &= \frac{1}{C_{\text{skin}}} \left[ R_{\downarrow} - \sigma_{\text{SB}} T_s^4 - \lambda \frac{T_s - T_{\text{deep}}}{d_{\text{soil}}} - \rho c_p C_H \sqrt{e + \delta} \cdot (T_s - T_a) \right] \label{eq:slow_ts}
\end{align}

\subsection{Parameterization of Closures and Functions}
\label{subsec:closures}

To close the system mathematically, the abstract couplings are mapped to physical planetary boundary layer functions:

\subsubsection{The Stratification Function $G(T_s)$}
The stable stratification is driven by the vertical temperature inversion between the reference atmospheric layer ($T_a$) and the ground skin temperature ($T_s$). It is parameterized using a smooth, continuous exponential activation function to ensure it is $C^\infty$ everywhere:
\begin{equation}
G(T_s) = \exp\left( \beta \frac{T_a - T_s}{T_a} \right) - 1
\label{eq:stratification_func}
\end{equation}

\begin{description}
    \item[Behavior:] When $T_s = T_a$ (neutral stratification), $G(T_s) = 0$, and buoyant destruction vanishes. When $T_s < T_a$ (nocturnal cooling), $G(T_s) > 0$, rapidly accelerating the buoyant destruction of TKE.
\end{description}

\subsubsection{The Turbulent Diffusivity/Drag Closure}
The momentum drag coefficient $\gamma$ and the surface thermodynamic sensible heat exchange coefficient $C_H$ are scaled with logarithmic roughness laws:
\begin{equation}
\gamma = \gamma_{\mathrm{eff}}\frac{\kappa^2}{\ln(h/z_{0m})^2} \cdot \frac{1}{h} \qquad C_H = \frac{\kappa^2}{\ln(h/z_{0m})\ln(h/z_{0h})}
\label{eq:drag_closures}
\end{equation}
Where $\kappa$ is the von K\'arm\'an constant, $h$ is the bulk boundary layer height scale, and $z_{0m}, z_{0h}$ are the momentum and thermal roughness lengths.

\subsection{Physical Parameters and Baseline Scaling}
\label{subsec:parameters}

To run this model numerically, the baseline parameter configuration scaled for a mid-latitude nocturnal boundary layer matching experimental CASES99 conditions is detailed in Table~\ref{tab:baseline_parameters}.

\begin{table}[ht!]
\centering
\caption{Baseline Physical Parameter Specifications for the 4D GSPT-SBL System}
\label{tab:baseline_parameters}
\begin{tabular}{llll}
\toprule
\textbf{Parameter} & \textbf{Description} & \textbf{Baseline Value} & \textbf{Units} \\ \midrule
\multicolumn{4}{l}{\textit{Environmental Controls and Geostrophic Forcing}} \\
$U_g$ & Zonal Geostrophic Wind & $10.0$ & $\text{m s}^{-1}$ \\
$V_g$ & Meridional Geostrophic Wind & $0.0$ & $\text{m s}^{-1}$ \\
$T_a$ & Reference Air Temperature ($z = h$) & $285.15$ & $\text{K}\ (12^\circ\text{C})$ \\
$T_{\text{deep}}$ & Deep Soil Core Temperature & $283.15$ & $\text{K}\ (10^\circ\text{C})$ \\
$R_{\downarrow}$ & Downward Longwave Radiation & $260.0$ & $\text{W m}^{-2}$ \\
$f$ & Coriolis Parameter ($\sim 45^\circ\text{N}$) & $1.0 \times 10^{-4}$ & $\text{s}^{-1}$ \\ \midrule
\multicolumn{4}{l}{\textit{Turbulence and Boundary Layer Internals}} \\
$\varepsilon$ & Singular Perturbation Parameter & $0.01$ & Dimensionless \\
$\delta$ & Background Mixing Parameter & $1.0 \times 10^{-4}$ & $\text{m}^2\text{s}^{-2}$ \\
$K$ & Buoyant Destruction Weight & $0.32$ & Dimensionless gain \\
$\beta$ & Stratification Sensitivity Parameter & $15.0$ & Dimensionless \\
$h$ & Coupling Bulk Scale Height & $50.0$ & $\text{m}$ \\
$\gamma_{\mathrm{eff}}$ & Drag Scaling Efficiency & $2.0$ & Dimensionless \\
$\eta$ & Shear Production Efficiency & $15.0$ & Dimensionless \\ \midrule
\multicolumn{4}{l}{\textit{Thermal Properties (Surface Energy Budget)}} \\
$\sigma_{\text{SB}}$ & Stefan-Boltzmann Constant & $5.67 \times 10^{-8}$ & $\text{W m}^{-2}\text{K}^{-4}$ \\
$\lambda$ & Soil Thermal Conductivity & $1.2$ & $\text{W m}^{-1}\text{K}^{-1}$ \\
$d_{\text{soil}}$ & Effective Soil Layer Depth & $0.5$ & $\text{m}$ \\
$\rho c_p$ & Volumetric Heat Capacity of Air & $1200.0$ & $\text{J m}^{-3}\text{K}^{-1}$ \\
$C_{\text{skin}}$ & Thermal Capacity of Skin Layer & $2.0 \times 10^{4}$ & $\text{J m}^{-2}\text{K}^{-1}$ \\ \bottomrule
\end{tabular}
\end{table}

\subsection{Geometric Features of the 4D Phase Space}
\label{subsec:geometry}

When analyzing this mathematical system, the geometric structures shift from previous lower-dimensional derivations in three important ways:

\begin{enumerate}
    \item \textbf{The Critical Manifold $\mathcal{M}_0$:} Defined strictly in the limit $\varepsilon \to 0$ by setting the fast equation to zero: $\mathcal{F}(e,U,V,T_s)=0$. On the active turbulent sheet, $e+\delta = l_0\Delta$ where $\Delta = \eta\,\gamma\,(U^2+V^2)-K\,G(T_s)$, so $e^* = l_0\Delta-\delta$.
    \item \textbf{The Trapped Laminar State:} If $\Delta \le 0$, no active turbulent root exists. The regularization parameter $\delta$ then holds the system near the background mixing floor, so $e$ remains close to the laminar branch.
    \item \textbf{The Fold Set ($\mathcal{C}_{\text{fold}}$):} The nonlinear surface energy budget (SEB) creates the parameter dependence required for the manifold to fold, but the geometric fold itself is defined by the simultaneous conditions $\mathcal{F}=0$ and $\partial_e\mathcal{F}=0$. Because $U^2 + V^2$ is circularly symmetric, the turbulent branch has paraboloid-of-revolution geometry in the $(U, V, e)$ coordinates for fixed $T_s$.
\end{enumerate}

\section{Bifurcation and Phase Space Analysis}
\label{sec:diagnostics}

Numerical integration of the fully coupled 4D system reveals a clear division between standard continuous equilibria and sudden dropout structures. Figure~\ref{fig:sbl_diagnostics} presents the trajectories tracked under the automated numerical engine, highlighting the geometric alignment of variables along the slow critical manifold.

\begin{figure}[ht!]
\centering
\includegraphics[width=0.95\linewidth]{reports/generated/figures/4d_sbl_diagnostics.png}
\caption{Trajectory trajectories and diagnostic parameters for the 4D GSPT-SBL system showing state transitions and the regularized floor interaction.}
\label{fig:sbl_diagnostics}
\end{figure}

\end{document}